% ================================================================
% DATOS EDITABLES DE LA TESIS
% Modifique este archivo; no es necesario tocar main.tex.
% ================================================================

\newcommand{\tituloTesis}{Título claro, específico y representativo de la tesis}
\newcommand{\autorTesis}{Nombre completo de la persona tesista}

\newcommand{\modalidadTitulacion}{Tesis}
\newcommand{\gradoObtener}{Ingeniero(a) de Software}

\newcommand{\gradoDirector}{Dr./Dra.}
\newcommand{\nombreDirector}{Nombre completo de la persona directora}

% Cambie a \tieneCodirectorfalse si la tesis no cuenta con codirección.
\tieneCodirectortrue
\newcommand{\gradoCodirector}{Dr./Dra.}
\newcommand{\nombreCodirector}{Nombre completo de la persona codirectora}

\newcommand{\gradoSinodalUno}{Dr./Dra.}
\newcommand{\nombreSinodalUno}{Nombre completo de la primera persona sinodal}
\newcommand{\gradoSinodalDos}{Dr./Dra.}
\newcommand{\nombreSinodalDos}{Nombre completo de la segunda persona sinodal}
\newcommand{\gradoSinodalTres}{Dr./Dra.}
\newcommand{\nombreSinodalTres}{Nombre completo de la tercera persona sinodal}

\newcommand{\lugarPublicacion}{Zacatecas, Zacatecas, México}
% Se actualiza al compilar (también en Overleaf). Si la fecha institucional
% requerida es distinta, sustituya temporalmente \mesAnioAutomatico por el
% texto correspondiente, por ejemplo: septiembre de 2026.
\newcommand{\mesAnioPublicacion}{\mesAnioAutomatico}

% Metadatos para el PDF; separe los términos con comas.
\newcommand{\palabrasClavePdf}{ingeniería de software, tesis, investigación}

% Mantenga las instrucciones y ejemplos visibles mientras redacta. Para generar
% la versión entregable cambie estas dos líneas a "false".
\mostrarInstruccionestrue
\mostrarEjemplostrue
